\section{Mathematical Prerequisites}
Quantum mechanics is formulated in the language of linear algebra. 
Understanding quantum algorithms requires fluency in the properties of complex vector spaces, 
inner products, and operator theory.

\subsection{Hilbert Spaces and Dirac Notation}
The state of an isolated quantum system is described by a unit vector in a complex Hilbert space, 
denoted $\mathcal{H}$. 

A Hilbert space is a vector space equipped with an inner product that is complete with respect to 
the norm induced by that inner product. For the finite-dimensional systems relevant to 
quantum computing (such as a system of $n$ qubits), the Hilbert space is isomorphic to 
$\mathbb{C}^{2^n}$.

We employ the Dirac \textit{bra-ket} notation standard in quantum physics:

\begin{itemize}
    \item \textbf{Ket} ($\ket{\psi}$): A column vector representing a quantum state.
    \item \textbf{Bra} ($\bra{\psi}$): The Hermitian conjugate (conjugate transpose) of the ket,
     representing a row vector in the dual space $\mathcal{H}^*$.
\end{itemize}

For a vector $\psi \in \mathbb{C}^d$, we write:
\begin{equation}
    \ket{\psi} = \begin{pmatrix} z_1 \\ z_2 \\ \vdots \\ z_d \end{pmatrix}, \quad 
    \bra{\psi} = \ket{\psi}^\dagger = \begin{pmatrix} z_1^* & z_2^* & \dots & z_d^* \end{pmatrix}
\end{equation}
where $z^*$ denotes the complex conjugate of $z$.

\subsection{Inner and Outer Products}
The \textbf{inner product} of two state vectors $\ket{\phi}$ and $\ket{\psi}$ is a scalar $c \in \mathbb{C}$, 
denoted by the bracket $\braket{\phi}{\psi}$. It satisfies the following properties:
\begin{enumerate}
    \item Conjugate symmetry: $\braket{\phi}{\psi} = \braket{\psi}{\phi}^*$.
    \item Linearity in the second argument: $\bra{\phi} (a\ket{\psi_1} + b\ket{\psi_2}) = a\braket{\phi}{\psi_1} + b\braket{\phi}{\psi_2}$.
    \item Positive definiteness: $\braket{\psi}{\psi} \geq 0$, with equality if and only if $\ket{\psi} = 0$.
\end{enumerate}
Two vectors are \textit{orthogonal} if their inner product is zero: $\braket{\phi}{\psi} = 0$. 
The norm of a vector is defined as $||\psi|| = \sqrt{\braket{\psi}{\psi}}$. 
To represent a valid physical quantum state, a state vector must be normalized to unity ($||\psi|| = 1$).

The \textbf{outer product} of $\ket{\psi}$ and $\ket{\phi}$ is an operator (matrix) $A = \ket{\psi}\bra{\phi}$. 
Applying this operator to a vector $\ket{\xi}$ yields a new vector proportional to $\ket{\psi}$:
\begin{equation}
    (\ket{\psi}\bra{\phi})\ket{\xi} = \ket{\psi}(\braket{\phi}{\xi}) = c \ket{\psi}
\end{equation}
where $c = \braket{\phi}{\xi}$ is a scalar.

\subsection{Linear Operators}
Quantum evolution is described by linear operators acting on $\mathcal{H}$. 
An operator $A$ is linear if $A(c_1\ket{\psi_1} + c_2\ket{\psi_2}) = c_1 A\ket{\psi_1} + c_2 A\ket{\psi_2}$.

\subsubsection{Hermitian Operators}
An operator $A$ is \textit{Hermitian} (or self-adjoint) if $A = A^\dagger$, 
where $A^\dagger$ is the conjugate transpose of $A$. 

Hermitian operators are crucial because they represent physical observables (such as energy or spin). 
The Spectral Theorem states that for any Hermitian operator $M$ on a finite-dimensional space, 
there exists an orthonormal basis of eigenvectors $\{\ket{m_i}\}$ with real eigenvalues $m_i$ such that:
\begin{equation}
    M = \sum_i m_i \ket{m_i}\bra{m_i}
\end{equation}
This decomposition implies that the outcome of a measurement of observable $M$ is always one of the real values $m_i$.

\subsubsection{Unitary Operators}
An operator $U$ is \textit{unitary} if its inverse is its adjoint:
\begin{equation}
    U^\dagger U = U U^\dagger = I
\end{equation}
Unitary operators preserve the inner product: $\bra{\psi} U^\dagger U \ket{\phi} = \braket{\psi}{\phi}$. 
Consequently, they preserve the norm of state vectors. This property is essential for quantum computing, 
as it ensures that the sum of probabilities of all possible outcomes remains $1$ during the evolution of the system. 

All closed quantum systems evolve via unitary transformations.